\section{Introduction}\label{sec:introduction}
A data owner may need to demonstrate that public cluster centers give a reported assignment cost on private records. Examples include an analyst checking a deployed model on sensitive embeddings and an auditor requesting costs for several proposed model updates. Zero-knowledge proofs establish the specified computation without disclosing its private witness \cite{gmr1991knowledge,groth2016}. The difficulty is repeated high-dimensional arithmetic: a direct nearest-center audit evaluates distances in $d$ dimensions for every record and every candidate center. In the completed 1,024-record MNIST control, our exact public-center Groth16 circuit requires 1.43 million constraints for a batch of eight 784-dimensional records.

The same committed records can be queried many times. This motivates a two-stage construction. Stage~1 proves once that compact statistics were derived correctly from every committed original row. Stage~2 authenticates the requested statistics and proves nearest-center assignments and their aggregate cost. The central systems question is which statistics allow useful later queries while avoiding repeated original-space work.

For one fixed center set, the geometry is elementary. An orthonormal basis spanning the center differences preserves all assignment comparisons; adding each record's squared residual norm recovers every center distance. This follows from Pythagoras and belongs to the established projection-plus-offset approach to clustering sketches \cite{feldman2013tiny,feldman2020tiny,cohen2015dimensionality}. We use this identity as an interface to a proof system. We do not propose a new dimensionality-reduction identity or a new coreset theorem.

Fixed-center answers alone are insufficient to justify a reusable representation. A stronger task-specific baseline computes and proves the exact winning center and minimum cost in Stage~1, then authenticates those two values in later subset queries. It is cheaper for a single fixed center set. Exact affine scores provide another compact baseline with the same fixed-center information. Both measured implementations specialize preprocessing to the selected centers, so their Stage~1 work is repeated when those centers change.

We therefore study audits under changing public centers. An analyst can compare proposed prototype updates against the same sensitive embedding catalog without disclosing individual records or assignments to an auditor. We augment the deployed-center span with principal components of training residuals, and store the resulting head coordinates and residual energy. New centers enter Stage~2 through public head coordinates and residual offsets. They change the query inputs without changing the catalog or the query setup. This yields an exact real-arithmetic distance when the new centers remain in the represented affine span, and an approximation otherwise. The approximation drops a residual cross term. Proving this approximate execution alone does not establish original-space accuracy. We add a circuit-checked error envelope that accounts for both the omitted cross term and the nonorthogonality of the rounded basis, yielding per-set cost intervals and certified assignment counts.

An exact method can also reuse setup: put the full centers in the public statement and compute original-space distances directly. We implement this public-center Full circuit and prove complete catalogs with it. The earlier 6.82--7.76-fold savings over center-specific caches are retained only as a historical implementation comparison. The completed serial campaign measures matched plain reuse, certified reuse and public-center Full, including an actual 1,024-record catalog and predeclared higher-rank controls. At 1,024 records, learned plain and certified reuse have 5.17- and 5.08-fold lower measured component cost than exact public-center Full. These alternatives provide different fidelity guarantees: the certificate can force an auditor to refuse an inexpensive approximate answer.

The contribution has three parts:
\begin{enumerate}[leftmargin=*]
\item An original-data-bound integer audit interface, with complete Stage~1 coverage, public-center reuse and independent public-only verification against approved tasks.
\item An optional rounded-basis accuracy relation: private row statistics yield a public lower bound on correct original-space assignments and an interval for exact quantized cost. An explicit policy can refuse a valid approximate proof when these bounds are insufficient.
\item A completed matched evaluation against exact public-center Full and equally reusable PCA/energy, including actual 1,024-record proofs, measured certificate overhead and higher-rank proof controls. All 6,848 campaign proofs pass public-only replay; historical timing cohorts and native budget grids remain separate.
\end{enumerate}

\section{Related work}\label{sec:related}
\paragraph{Projection and residual offsets}
Feldman, Schmidt and Sohler reduce clustering dimension while retaining an additive cost independent of the queried solution \cite{feldman2013tiny,feldman2020tiny}. Their arXiv extended version states the matrix Pythagorean identity (Claim~3) and, for centers in a subspace, the decomposition into projection cost and orthogonal residual cost (Eq.~9). Cohen et al. formalize projection-cost-preserving sketches with a solution-independent additive constant \cite{cohen2015dimensionality}. These results address approximation over classes of clustering solutions. Here we keep one energy value per authenticated record so that arbitrary requested subsets can be audited, and prove the deterministic derivation from the original commitment. The fixed-center identity is a special case of this classical geometry; the reuse circuit and its measured trade-offs are the focus.

\paragraph{Verifiable clustering and private analytics}
Prior work already combines private $k$-means and zero-knowledge techniques. Duan and Canny use random projections in zero-knowledge checks of vector norm bounds for private distributed data mining \cite{duan2008private}. Patel, Patel and Jinwala use secret sharing with a code-based identification scheme for horizontally partitioned clustering against malicious participants \cite{patel2013clustering}. Zhang, Li and Chang describe verifiable outsourced privacy-preserving $k$-means using secure distance and minimum computations \cite{zhang2021verifiable}. Repetto et al. prove PAC-private mechanisms, including $k$-means, using RISC-Zero \cite{repetto2026pac}. These works address input validity, distributed computation or outsourced clustering, whereas our task audits successive public center sets against an unchanged committed catalog. In nearby vector search, V3DB proves execution of published approximate-search semantics over committed snapshots \cite{v3db2026}. Our relation instead proves exhaustive nearest-center comparisons on authenticated statistics and checks an envelope for original-space quantized answers. Neither projection-based ZK checks, ZK clustering nor audits over committed vector data are new categories.

\paragraph{Proof-oriented representations}
ezDPS studies a ZK inference pipeline including PCA, and Artemis develops commit-and-prove techniques for external data or model commitments \cite{ezdps2023,artemis2024}. Our PCA baseline proves an approved projected computation in the same backend; it does not reproduce ezDPS or prove PCA training. We use Groth16 over BN254 and Poseidon2 row-Merkle authentication \cite{groth2016,poseidon22023}. The contribution is a representation and circuit design within this setting, not a new proof primitive.

\paragraph{Privacy-preserving prototype models}
Schubert et al. implement a nearest-prototype learning model using fully homomorphic encryption and study the effects of integer encoding and distance computation \cite{schubert2026prototype}. Their objective is computation over encrypted data. Our objective is a zero-knowledge audit of specified nearest-center answers over committed records. Both settings make representation cost and integer semantics important, but encryption, proof soundness and approximation fidelity address different guarantees. The present experiments audit public centers; they do not prove private model training or neural feature extraction.

\paragraph{Data-oblivious alternatives}
JL and sparse JL preserve geometry for an independently fixed point set under their stated distribution and dimension assumptions \cite{johnson1984extensions,achlioptas2003database,kane2014sparser,makarychev2019performance}. A post-commit random map avoids requiring training data to choose a projection. We retain that protocol option and evaluate sparse-JL estimators, but the tested tails do not outperform PCA with energy on matched new-center cells. The binding guarantee for the learned representation comes from Stage~1 proofs, independently of whether a random tail is present.
